% GENERATED FILE. Edit canonical lessons, then run npm run generate:books.
% canonical-source-hash: fnv1a64:1ef9431dc38031f2
% canonical-lessons: KA-C377-tarka, KA-C377-samarthisu, KA-C377-vivada, KA-C377-aksepisu, KA-C377-samanjasa

\chapter{Logic and Objections}
\label{ch:ka-logic-and-objections}
\hlchaptermodality{377}

\begin{chapteropening}
\textbf{By the end of this chapter:} \emph{I can say logic, to justify, a controversy, to object and reasonable.}

Complete the last lesson of chapter 377: I can say logic, to justify, a controversy, to object and reasonable.
\end{chapteropening}

\section[\texorpdfstring{\kn{ತರ್ಕ} (tarka)}{ತರ್ಕ (tarka)}]{\kn{ತರ್ಕ} (tarka) — logic; reasoning}

\label{lesson:KA-C377-tarka}



\begin{quote}
Before the new one: say the Kannada for importance, then the Kannada for to base on.
\end{quote}

\subsection*{You'll want to know: \kn{ತರ್ಕ}}
\textbf{\kn{ತರ್ಕ}} — \emph{tarka} — \textquotedblleft{}logic; reasoning\textquotedblright{}.

\begin{quote}
\kn{ಅವನ} \kn{ತರ್ಕ} \kn{ನನಗೆ} \kn{ಸರಿ} \kn{ಅನಿಸಲಿಲ್ಲ}.
\end{quote}

\begin{quote}
\emph{avana tarka nanage sari anisalilla.}
\end{quote}

\begin{quote}
His reasoning did not seem right to me.
\end{quote}

\kn{ತರ್ಕಬದ್ಧ} means reasoned, logical.

\begin{cousinweb}
\textbf{\kn{ತರ್ಕ}} (\emph{tarka}) is borrowed from Sanskrit \emph{tarka}, \textquotedblleft{}conjecture, reasoning\textquotedblright{} (CDIAL 5714).
\end{cousinweb}

\begin{grammarlens}[title={What you've built}]
One more word for giving reasons and opinions.
\end{grammarlens}

\subsection*{Your turn}
\begin{itemize}
\raggedright
  \item \emph{Say it:} \emph{tarka}
  \item \emph{Say it:} \emph{tarka}, once more
  \item \emph{Say it:} \emph{avana tarka nanage sari anisalilla.} — His reasoning did not seem right to me.
  \item \emph{Recall:} say the Kannada for importance, then the Kannada for to base on, then say \emph{tarka} again
\end{itemize}

\begin{tcolorbox}[breakable,colback=teal!4,colframe=teal!35!black,title={Before you move on}]
What does \kn{ತರ್ಕ} mean? (\textquotedblleft{}Logic; reasoning\textquotedblright{}.) Say it once more.
\end{tcolorbox}
\section[\texorpdfstring{\kn{ಸಮರ್ಥಿಸು} (samarthisu)}{ಸಮರ್ಥಿಸು (samarthisu)}]{\kn{ಸಮರ್ಥಿಸು} (samarthisu) — to justify; to defend}

\label{lesson:KA-C377-samarthisu}



\begin{quote}
Before the new one: say the Kannada for to base on, then the Kannada for logic.
\end{quote}

\subsection*{You'll want to know: \kn{ಸಮರ್ಥಿಸು}}
\textbf{\kn{ಸಮರ್ಥಿಸು}} — \emph{samarthisu} — \textquotedblleft{}to justify; to defend\textquotedblright{}.

\begin{quote}
\kn{ಅವಳು} \kn{ತನ್ನ} \kn{ನಿರ್ಧಾರವನ್ನು} \kn{ಸಮರ್ಥಿಸಿದಳು}.
\end{quote}

\begin{quote}
\emph{avaḷu tanna nirdhāravannu samarthisidaḷu.}
\end{quote}

\begin{quote}
She defended her decision.
\end{quote}

Past (I …): \kn{ಸಮರ್ಥಿಸಿದೆ}, \emph{samarthiside}.

What you defend takes -\kn{ಅನ್ನು}. \kn{ಸಮರ್ಥಿಸಿಕೊ}, to defend oneself.

\begin{cousinweb}
\textbf{\kn{ಸಮರ್ಥಿಸು}} (\emph{samarthisu}) is built on Sanskrit \emph{samartha}, \textquotedblleft{}fit, able, forcible\textquotedblright{} (CDIAL 13189), with the Kannada verb ending -\emph{isu}: to show that something is fit.
\end{cousinweb}

\begin{grammarlens}[title={What you've built}]
One more word for giving reasons and opinions.
\end{grammarlens}

\subsection*{Your turn}
\begin{itemize}
\raggedright
  \item \emph{Say it:} \emph{samarthisu}
  \item \emph{Say it:} \emph{samarthisu}, once more
  \item \emph{Say it:} \emph{avaḷu tanna nirdhāravannu samarthisidaḷu.} — She defended her decision.
  \item \emph{Recall:} say the Kannada for to base on, then the Kannada for logic, then say \emph{samarthisu} again
\end{itemize}

\begin{tcolorbox}[breakable,colback=teal!4,colframe=teal!35!black,title={Before you move on}]
What does \kn{ಸಮರ್ಥಿಸು} mean? (\textquotedblleft{}To justify; to defend\textquotedblright{}.) Say it once more.
\end{tcolorbox}
\section[\texorpdfstring{\kn{ವಿವಾದ} (vivāda)}{ವಿವಾದ (vivāda)}]{\kn{ವಿವಾದ} (vivāda) — a controversy; a dispute}

\label{lesson:KA-C377-vivada}



\begin{quote}
Before the new one: say the Kannada for logic, then the Kannada for to justify.
\end{quote}

\subsection*{You'll want to know: \kn{ವಿವಾದ}}
\textbf{\kn{ವಿವಾದ}} — \emph{vivāda} — \textquotedblleft{}a controversy; a dispute\textquotedblright{}.

\begin{quote}
\kn{ಹೊಸ} \kn{ನಿಯಮದ} \kn{ಬಗ್ಗೆ} \kn{ದೊಡ್ಡ} \kn{ವಿವಾದ} \kn{ಆಯಿತು}.
\end{quote}

\begin{quote}
\emph{hosa niyamada bagge doḍḍa vivāda āyitu.}
\end{quote}

\begin{quote}
There was a big controversy about the new rule.
\end{quote}

\kn{ವಿವಾದಾತ್ಮಕ}, controversial.

\begin{cousinweb}
\textbf{\kn{ವಿವಾದ}} (\emph{vivāda}) is borrowed from Sanskrit \emph{vivāda}, \textquotedblleft{}a dispute, a quarrel, a lawsuit\textquotedblright{} (Monier-Williams), from \emph{vi-} \textquotedblleft{}apart\textquotedblright{} and \emph{vad} \textquotedblleft{}to speak\textquotedblright{}.
\end{cousinweb}

\begin{grammarlens}[title={What you've built}]
One more word for giving reasons and opinions.
\end{grammarlens}

\subsection*{Your turn}
\begin{itemize}
\raggedright
  \item \emph{Say it:} \emph{vivāda}
  \item \emph{Say it:} \emph{vivāda}, once more
  \item \emph{Say it:} \emph{hosa niyamada bagge doḍḍa vivāda āyitu.} — There was a big controversy about the new rule.
  \item \emph{Recall:} say the Kannada for logic, then the Kannada for to justify, then say \emph{vivāda} again
\end{itemize}

\begin{tcolorbox}[breakable,colback=teal!4,colframe=teal!35!black,title={Before you move on}]
What does \kn{ವಿವಾದ} mean? (\textquotedblleft{}A controversy; a dispute\textquotedblright{}.) Say it once more.
\end{tcolorbox}
\section[\texorpdfstring{\kn{ಆಕ್ಷೇಪಿಸು} (ākṣēpisu)}{ಆಕ್ಷೇಪಿಸು (ākṣēpisu)}]{\kn{ಆಕ್ಷೇಪಿಸು} (ākṣēpisu) — to object}

\label{lesson:KA-C377-aksepisu}



\begin{quote}
Before the new one: say the Kannada for to justify, then the Kannada for a controversy.
\end{quote}

\subsection*{You'll want to know: \kn{ಆಕ್ಷೇಪಿಸು}}
\textbf{\kn{ಆಕ್ಷೇಪಿಸು}} — \emph{ākṣēpisu} — \textquotedblleft{}to object\textquotedblright{}.

\begin{quote}
\kn{ಊರಿನ} \kn{ಜನರು} \kn{ಹೊಸ} \kn{ಕಾರ್ಖಾನೆಗೆ} \kn{ಆಕ್ಷೇಪಿಸಿದರು}.
\end{quote}

\begin{quote}
\emph{ūrina janaru hosa kārkhānege ākṣēpisidaru.}
\end{quote}

\begin{quote}
The villagers objected to the new factory.
\end{quote}

Past (I …): \kn{ಆಕ್ಷೇಪಿಸಿದೆ}, \emph{ākṣēpiside}.

What you object to takes -\kn{ಗೆ}. \kn{ಆಕ್ಷೇಪ}, an objection.

\begin{cousinweb}
\textbf{\kn{ಆಕ್ಷೇಪಿಸು}} (\emph{ākṣēpisu}) is built on Sanskrit \emph{ākṣepa}, which among other senses means \textquotedblleft{}an objection\textquotedblright{} and \textquotedblleft{}a reproach\textquotedblright{} (Monier-Williams), with the verb ending -\emph{isu}.
\end{cousinweb}

\begin{grammarlens}[title={What you've built}]
One more word for giving reasons and opinions.
\end{grammarlens}

\subsection*{Your turn}
\begin{itemize}
\raggedright
  \item \emph{Say it:} \emph{ākṣēpisu}
  \item \emph{Say it:} \emph{ākṣēpisu}, once more
  \item \emph{Say it:} \emph{ūrina janaru hosa kārkhānege ākṣēpisidaru.} — The villagers objected to the new factory.
  \item \emph{Recall:} say the Kannada for to justify, then the Kannada for a controversy, then say \emph{ākṣēpisu} again
\end{itemize}

\begin{tcolorbox}[breakable,colback=teal!4,colframe=teal!35!black,title={Before you move on}]
What does \kn{ಆಕ್ಷೇಪಿಸು} mean? (\textquotedblleft{}To object\textquotedblright{}.) Say it once more.
\end{tcolorbox}
\section[\texorpdfstring{\kn{ಸಮಂಜಸ} (samañjasa)}{ಸಮಂಜಸ (samañjasa)}]{\kn{ಸಮಂಜಸ} (samañjasa) — reasonable; fitting}

\label{lesson:KA-C377-samanjasa}



\begin{quote}
Before the new one: say the Kannada for a controversy, then the Kannada for to object.
\end{quote}

\subsection*{You'll want to know: \kn{ಸಮಂಜಸ}}
\textbf{\kn{ಸಮಂಜಸ}} — \emph{samañjasa} — \textquotedblleft{}reasonable; fitting\textquotedblright{}.

\begin{quote}
\kn{ಅವರು} \kn{ಕೇಳಿದ} \kn{ಬೆಲೆ} \kn{ಸಮಂಜಸವಾಗಿತ್ತು}.
\end{quote}

\begin{quote}
\emph{avaru kēḷida bele samañjasavāgittu.}
\end{quote}

\begin{quote}
The price they asked was reasonable.
\end{quote}

Before a noun: \kn{ಸಮಂಜಸವಾದ} \kn{ಉತ್ತರ}, a reasonable answer.

\begin{cousinweb}
\textbf{\kn{ಸಮಂಜಸ}} (\emph{samañjasa}) is borrowed from Sanskrit \emph{samañjasa}, \textquotedblleft{}proper, right, sound\textquotedblright{} (Monier-Williams).
\end{cousinweb}

\begin{grammarlens}[title={What you've built}]
One more word for giving reasons and opinions.
\end{grammarlens}

\subsection*{Your turn}
\begin{itemize}
\raggedright
  \item \emph{Say it:} \emph{samañjasa}
  \item \emph{Say it:} \emph{samañjasa}, once more
  \item \emph{Say it:} \emph{avaru kēḷida bele samañjasavāgittu.} — The price they asked was reasonable.
  \item \emph{Recall:} say the Kannada for a controversy, then the Kannada for to object, then say \emph{samañjasa} again
\end{itemize}

\begin{tcolorbox}[breakable,colback=teal!4,colframe=teal!35!black,title={Before you move on}]
What does \kn{ಸಮಂಜಸ} mean? (\textquotedblleft{}Reasonable; fitting\textquotedblright{}.) Say it once more.
\end{tcolorbox}
